\documentclass[../../../include/open-logic-section]{subfiles}

\begin{document}

\olfileid{his}{set}{pathology}
\olsection{پدیده‌های نامتعارف}

اما روشن شد که تعریف \emph{حد} امکان ساختن نمونه‌هایی نسبتاً «نامتعارف»
را فراهم می‌کند.

بولتسانو در حدود دههٔ ۱۸۳۰ تابعی کشف کرد که \emph{همه‌جا پیوسته} اما
\emph{هیچ‌جا مشتق‌پذیر} نبود. (متأسفانه بولتسانو این نتیجه را هرگز منتشر
نکرد؛ ریاضی‌دانان نخست در ۱۸۷۲، به لطف کشف مستقل همان ایده به دست
وایرشتراس، با آن آشنا شدند.)\footnote{تاریخ این موضوع با جزئیاتی بسیار
کامل در پانوشت‌های مقالهٔ ویکی‌پدیا دربارهٔ
\href{http://en.wikipedia.org/wiki/Weierstrass_function}{تابع
وایرشتراس} مستند شده است.} این نتیجه، دست‌کم، بسیار شگفت‌آور بود. یافتن
توابعی مانند $|x|$ که همه‌جا پیوسته‌اند اما در نقطه‌ای خاص مشتق‌پذیر
نیستند آسان است. ولی تابعی که همه‌جا پیوسته و \emph{هیچ‌جا} مشتق‌پذیر
نباشد، موجودی کاملاً متفاوت است. لحظه‌ای بیندیشید که چگونه می‌توان چنین
تابعی را رسم کرد. برای تضمین پیوستگی باید بتوانید بی‌آنکه قلم را از روی
کاغذ بردارید آن را بکشید؛ اما برای آنکه هیچ‌جا مشتق‌پذیر نباشد، باید
پیوسته جهت حرکت قلم را ناگهان عوض کنید.

«پدیده‌های نامتعارف» دیگری نیز از پی آمدند. کانتور در ۵ ژانویهٔ ۱۸۷۴
نامه‌ای به ددکیند نوشت و این مسئله را مطرح کرد:
\begin{quote}
آیا می‌توان سطحی (مثلاً مربعی همراه با مرزش) را با خطی (مثلاً پاره‌خطی
همراه با دو سرش) تناظر یک‌به‌یک داد، به‌گونه‌ای که به هر نقطهٔ سطح نقطه‌ای
از خط متناظر باشد و، برعکس، به هر نقطهٔ خط نقطه‌ای از سطح؟

در این لحظه هنوز به نظرم می‌رسد که پاسخ این پرسش بسیار دشوار است---هرچند
در اینجا نیز آدمی چنان به گفتن \emph{نه} کشانده می‌شود که مایل است برهان
را تقریباً زائد بشمارد. [نقل‌شده در \citealt{Gouvea2011}]
\end{quote}
اما کانتور در ۱۸۷۷ ثابت کرد که اشتباه می‌کرده است. در واقع، یک خط و یک
مربع دقیقاً به یک تعداد نقطه دارند. او در ۲۹ ژوئن ۱۸۷۷ به ددکیند نوشت:
«\emph{je le vois, mais je ne le crois pas}»؛ یعنی «می‌بینمش، اما باورش
نمی‌کنم». در «روایت پذیرفته‌شده» از تاریخ ریاضیات، این سخن را اغلب نشانهٔ
آن می‌گیرند که این نتایج جدید تا چه اندازه برای ریاضی‌دانان آن روزگار
\emph{به معنای واقعی ناباورکردنی} بوده‌اند. (این مکاتبات در
\citet{Gouvea2011} آمده‌اند؛ در \olref[his][set][mythology]{sec} به آن
بازمی‌گردیم و طرح برهان کانتور در \olref[his][set][cantorplane]{sec}
ارائه می‌شود.)

پئانو با الهام از نتیجهٔ کانتور به این اندیشید که آیا می‌توان خطی را
\emph{به‌طور پیوسته} بر صفحه نگاشت. حاصل، \emph{منحنی پُرکنندهٔ فضا}
می‌بود. پئانو در \citeyear{Peano1890} دقیقاً چنین منحنی‌ای ساخت. این نتیجه
به‌راستی خلاف شهود است: اقلیدس خط را «طولی بی‌عرض» تعریف کرده بود (کتاب
اول، تعریف ۲)، اما پئانو نشان داد که با پیچاندن مناسب خط می‌توان طول آن را
به عرض بدل کرد. هیلبرت در \citeyear{Hilbert1891} منحنی پُرکنندهٔ فضای
اندکی شهودی‌تری را همراه با چند تصویر توضیح داد. منحنی مرحله‌به‌مرحله
ساخته می‌شود و در زیر شش مرحلهٔ نخست آن آمده است:
\begin{center}
\begin{tikzpicture}
\begin{scope}[xshift=20pt, yshift=20pt]
	\draw [oldiagcolorC, l-system={Hilbert curve, step=40pt, angle=90, axiom=L, order=1}] lindenmayer system;
\end{scope}
\begin{scope}[xshift=110pt, yshift=10pt]
	\draw [oldiagcolorC, l-system={Hilbert curve, step=20pt, angle=90, axiom=L, order=2}] lindenmayer system;
\end{scope}
\begin{scope}[xshift=205pt, yshift=5pt]
	\draw [oldiagcolorC, l-system={Hilbert curve, step=10pt, angle=90, axiom=L, order=3}] lindenmayer system;
\end{scope}
\begin{scope}[xshift=2.5pt, yshift=-97.5pt]
	\draw [oldiagcolorC, l-system={Hilbert curve, step=5pt, angle=90, axiom=L, order=4}] lindenmayer system;
\end{scope}
\begin{scope}[xshift=101.25pt, yshift=-98.75pt]
	\draw [oldiagcolorC, l-system={Hilbert curve, step=2.5pt, angle=90, axiom=L, order=5}] lindenmayer system;
\end{scope}
\begin{scope}[xshift=200.625pt, yshift=-99.375pt]
	\draw [oldiagcolorC, l-system={Hilbert curve, step=1.25pt, angle=90, axiom=L, order=6}] lindenmayer system;
\end{scope}
\draw[gray] (0pt,0pt) rectangle (80pt, 80pt);
\draw[gray] (100pt,0pt) rectangle (180pt, 80pt);
\draw[gray] (200pt,0pt) rectangle (280pt, 80pt);
\draw[gray] (0pt,-100pt) rectangle (80pt, -20pt);
\draw[gray] (100pt,-100pt) rectangle (180pt, -20pt);
\draw[gray] (200pt,-100pt) rectangle (280pt, -20pt);
\end{tikzpicture}
\end{center}
در حد---مفهومی که تا آن زمان تعریفی دقیق یافته بود---تمام مربع یکپارچه
با رنگ !!{colorC} پُر می‌شود. در ضمن، منحنی هیلبرت همه‌جا پیوسته اما
هیچ‌جا مشتق‌پذیر نیست، و دلیل شهودی آن این است که تابع در حد نامتناهی در
هر لحظه جهت خود را ناگهان تغییر می‌دهد. (ساخت هیلبرت را با جزئیات بیشتر
در \olref[his][set][hilbertcurve]{sec} طرح خواهیم کرد.)

خوب یا بد، این ساخت‌های هندسی «نامتعارف» دلیلی برای تردید در توسل به شهود
هندسی تلقی شدند. آن‌ها تقریباً به \emph{تبلیغاتی} برای شیوه‌ای نو در
ریاضی‌ورزی بدل شدند که سرانجام به نظریهٔ مجموعه‌ها می‌انجامید. در
اسطوره‌سازی‌های بعدی این رشته بارها تکرار شد که این نتایج هم کاملاً دقیق
و هم کاملاً تکان‌دهنده بودند. ازاین‌رو دو مقصود را برآورده می‌کردند:
هشداری دربارهٔ تکیه بر شهود هندسی و نمایشی از باروری شیوه‌های نوین تفکر.

\end{document}
